\documentclass[12pt]{article}
\usepackage[utf8]{inputenc}
\usepackage[T1]{fontenc}
\usepackage[margin=1in]{geometry}
\usepackage{setspace}
\usepackage{xcolor}
\usepackage[colorlinks=true, urlcolor=blue!60!black, linkcolor=blue!60!black]{hyperref}
\setlength{\parindent}{0em}
\setlength{\parskip}{0.6em}
\renewcommand{\baselinestretch}{1.0}
\begin{document}
\begin{center}
{\LARGE \textit{Plyler v. Doe}, 457 U.S. 202 (1982)\par}
\vspace{0.5em}
{\large Case Brief --- Two-Page Edition\par}
\end{center}

\section{Memory Jogger}
Texas could not deny free public schooling to undocumented children; a law imposing a lifetime hardship on innocent children must further a substantial state goal (\href{https://supreme.justia.com/cases/federal/us/457/202/}{457 U.S. 202}).

\section{Detailed Case Facts}
A 1975 Texas law withheld state funds for educating children not ``legally admitted'' into the United States and let districts exclude them; Tyler ISD charged them \$1,000 a year in tuition.

\section{Procedural History}
District courts enjoined the law after finding it did not deter immigration or improve schools; the Fifth Circuit affirmed in both consolidated cases.

\section{Judicial Votes}
Affirmed 5--4. Brennan wrote for Marshall, Blackmun, Powell, and Stevens; Marshall, Blackmun, and Powell also concurred separately; Burger, with White, Rehnquist, and O'Connor, dissented.

\section{Holding}
Undocumented immigrants are ``persons within its jurisdiction.'' Though they are not a suspect class and education is not a fundamental right, excluding their children fails unless it furthers a substantial state goal.

\section{Analysis of Opinions}
\textbf{Majority:} Children ``can affect neither their parents' conduct nor their own status,'' and exclusion risks a permanent ``shadow population''; education sustains democratic institutions. Tuition was a ``ludicrously ineffectual attempt to stem the tide of illegal immigration.'' \textbf{Blackmun and Powell:} total denial is like denying the vote; illegitimacy cases support heightened review. \textbf{Burger:} the Court acted as ``Platonic Guardians''; rational basis should control.

\section{Examples: Future Applications}
\textbf{Same-Side:} requiring proof of lawful status to enroll is invalid. \textbf{Opposite-Side:} a neutral residency rule is valid (\textit{Martinez v. Bynum}). \textbf{Fence-Sitter:} requiring schools to report students' immigration status.

\section{Critique}
Critics say \textit{Plyler} combined two non-triggers into ad hoc scrutiny better suited to preemption. \textit{Rodriguez} and \textit{Kadrmas} mark its limits.

\section{Key Quotations}
\begin{quote}
``Section 21.031 imposes a lifetime hardship on a discrete class of children not accountable for their disabling status. The stigma of illiteracy will mark them for the rest of their lives. By denying these children a basic education, we deny them the ability to live within the structure of our civic institutions, and foreclose any realistic possibility that they will contribute in even the smallest way to the progress of our Nation. In determining the rationality of \S~21.031, we may appropriately take into account its costs to the Nation and to the innocent children who are its victims. In light of these countervailing costs, the discrimination contained in \S~21.031 can hardly be considered rational unless it furthers some substantial goal of the State.''
\end{quote}
\noindent --- Justice Brennan, Opinion of the Court, 457 U.S. at 223--24.

\end{document}
